% GENERATED FILE. Edit canonical lessons, then run npm run generate:books.
% canonical-source-hash: fnv1a64:8d6140ddd4994651
% canonical-lessons: JA-C182-harau, JA-C182-otsuri, JA-C182-genkin, JA-C182-satsu, JA-C182-kozeni

\chapter{Paying, Cash and Coins}
\label{ch:ja-paying-cash-and-coins}
\hlchaptermodality{182}

\begin{chapteropening}
\textbf{By the end of this chapter:} \emph{I can pay and talk about cash and change.}

Complete the last lesson of chapter 182: I can pay and talk about cash and change.
\end{chapteropening}

\section[\texorpdfstring{\ja{はらう} (harau)}{はらう (harau)}]{\ja{はらう} (harau) — to pay}

\label{lesson:JA-C182-harau}



\begin{quote}
Before the new one: say the Japanese for to sell, then the Japanese for to sell well.
\end{quote}

\subsection*{You'll want to know: \ja{はらう}}
\textbf{\ja{はらう}} — \emph{harau} — \textquotedblleft{}to pay\textquotedblright{}.

\begin{quote}
\textbf{\ja{ここで} \ja{おかねを} \ja{はらいます。}} — \emph{koko de okane o haraimasu} — \textquotedblleft{}Customers pay here.\textquotedblright{}
\end{quote}

Polite \textbf{\ja{はらいます}}, \emph{haraimasu}; a godan verb.

\begin{cousinweb}
\textbf{\ja{はらう}} (\emph{harau}) is a native Japanese word (wago), inherited rather than borrowed from Chinese; its kanji was attached later as a kun'yomi (JMdict, KANJIDIC2).
\end{cousinweb}

\begin{grammarlens}[title={What you've built}]
One more word for everyday life at A2.
\end{grammarlens}

\subsection*{Your turn}
\begin{itemize}
\raggedright
  \item \emph{Say it:} \emph{harau}
  \item \emph{Say it:} \emph{harau}, once more
  \item \emph{Say it:} \emph{koko de okane o haraimasu} — Customers pay here.
  \item \emph{Recall:} say the Japanese for to sell, then the Japanese for to sell well, then say \emph{harau} again
\end{itemize}

\begin{tcolorbox}[breakable,colback=teal!4,colframe=teal!35!black,title={Before you move on}]
What does \ja{はらう} mean? (\textquotedblleft{}To pay\textquotedblright{}.) Say it once more.
\end{tcolorbox}
\section[\texorpdfstring{\ja{おつり} (otsuri)}{おつり (otsuri)}]{\ja{おつり} (otsuri) — change (money back)}

\label{lesson:JA-C182-otsuri}



\begin{quote}
Before the new one: say the Japanese for to sell well, then the Japanese for to pay.
\end{quote}

\subsection*{You'll want to know: \ja{おつり}}
\textbf{\ja{おつり}} — \emph{otsuri} — \textquotedblleft{}change (money back)\textquotedblright{}.

\begin{quote}
\textbf{\ja{おつりは} \ja{じゅうえんです。}} — \emph{otsuri wa jūen desu} — \textquotedblleft{}The change is ten yen.\textquotedblright{}
\end{quote}

\begin{cousinweb}
\textbf{\ja{おつり}} (\emph{otsuri}) is built from polite \emph{o-} and native \emph{tsuri}, written with the same kanji as \emph{tsuri}, \textquotedblleft{}fishing\textquotedblright{}; how the money sense arose is not settled (JMdict, KANJIDIC2).
\end{cousinweb}

\begin{grammarlens}[title={What you've built}]
One more word for everyday life at A2.
\end{grammarlens}

\subsection*{Your turn}
\begin{itemize}
\raggedright
  \item \emph{Say it:} \emph{otsuri}
  \item \emph{Say it:} \emph{otsuri}, once more
  \item \emph{Say it:} \emph{otsuri wa jūen desu} — The change is ten yen.
  \item \emph{Recall:} say the Japanese for to sell well, then the Japanese for to pay, then say \emph{otsuri} again
\end{itemize}

\begin{tcolorbox}[breakable,colback=teal!4,colframe=teal!35!black,title={Before you move on}]
What does \ja{おつり} mean? (\textquotedblleft{}Change (money back)\textquotedblright{}.) Say it once more.
\end{tcolorbox}
\section[\texorpdfstring{\ja{げんきん} (genkin)}{げんきん (genkin)}]{\ja{げんきん} (genkin) — cash}

\label{lesson:JA-C182-genkin}



\begin{quote}
Before the new one: say the Japanese for to pay, then the Japanese for change.
\end{quote}

\subsection*{You'll want to know: \ja{げんきん}}
\textbf{\ja{げんきん}} — \emph{genkin} — \textquotedblleft{}cash\textquotedblright{}.

\begin{quote}
\textbf{\ja{げんきんで} \ja{はらいます。}} — \emph{genkin de haraimasu} — \textquotedblleft{}I pay in cash.\textquotedblright{}
\end{quote}

\begin{cousinweb}
\textbf{\ja{げんきん}} (\emph{genkin}) is Sino-Japanese (kango): two Chinese characters read with their on'yomi, readings borrowed from Chinese. They are \emph{gen}, \textquotedblleft{}present, actual\textquotedblright{} (Middle Chinese \emph{henH}) and \emph{kin}, \textquotedblleft{}gold, money\textquotedblright{} (\emph{kim}) (KANJIDIC2; Middle Chinese in Baxter's transcription).
\end{cousinweb}

\begin{grammarlens}[title={What you've built}]
One more word for everyday life at A2.
\end{grammarlens}

\subsection*{Your turn}
\begin{itemize}
\raggedright
  \item \emph{Say it:} \emph{genkin}
  \item \emph{Say it:} \emph{genkin}, once more
  \item \emph{Say it:} \emph{genkin de haraimasu} — I pay in cash.
  \item \emph{Recall:} say the Japanese for to pay, then the Japanese for change, then say \emph{genkin} again
\end{itemize}

\begin{tcolorbox}[breakable,colback=teal!4,colframe=teal!35!black,title={Before you move on}]
What does \ja{げんきん} mean? (\textquotedblleft{}Cash\textquotedblright{}.) Say it once more.
\end{tcolorbox}
\section[\texorpdfstring{\ja{さつ} (satsu)}{さつ (satsu)}]{\ja{さつ} (satsu) — banknote; paper money}

\label{lesson:JA-C182-satsu}



\begin{quote}
Before the new one: say the Japanese for change, then the Japanese for cash.
\end{quote}

\subsection*{You'll want to know: \ja{さつ}}
\textbf{\ja{さつ}} — \emph{satsu} — \textquotedblleft{}banknote; paper money\textquotedblright{}.

\begin{quote}
\textbf{\ja{さいふに} \ja{さつが} \ja{ありません。}} — \emph{saifu ni satsu ga arimasen} — \textquotedblleft{}There are no banknotes in my wallet.\textquotedblright{}
\end{quote}

\begin{cousinweb}
\textbf{\ja{さつ}} (\emph{satsu}) is Sino-Japanese (kango): one Chinese character, \textquotedblleft{}tag, paper money\textquotedblright{}, read with its on'yomi, a reading borrowed from Chinese (Middle Chinese \emph{tsreat}) (KANJIDIC2; Middle Chinese in Baxter's transcription).
\end{cousinweb}

\begin{grammarlens}[title={What you've built}]
One more word for everyday life at A2.
\end{grammarlens}

\subsection*{Your turn}
\begin{itemize}
\raggedright
  \item \emph{Say it:} \emph{satsu}
  \item \emph{Say it:} \emph{satsu}, once more
  \item \emph{Say it:} \emph{saifu ni satsu ga arimasen} — There are no banknotes in my wallet.
  \item \emph{Recall:} say the Japanese for change, then the Japanese for cash, then say \emph{satsu} again
\end{itemize}

\begin{tcolorbox}[breakable,colback=teal!4,colframe=teal!35!black,title={Before you move on}]
What does \ja{さつ} mean? (\textquotedblleft{}Banknote; paper money\textquotedblright{}.) Say it once more.
\end{tcolorbox}
\section[\texorpdfstring{\ja{こぜに} (kozeni)}{こぜに (kozeni)}]{\ja{こぜに} (kozeni) — small coins; small change}

\label{lesson:JA-C182-kozeni}



\begin{quote}
Before the new one: say the Japanese for cash, then the Japanese for banknote.
\end{quote}

\subsection*{You'll want to know: \ja{こぜに}}
\textbf{\ja{こぜに}} — \emph{kozeni} — \textquotedblleft{}small coins; small change\textquotedblright{}.

\begin{quote}
\textbf{\ja{こぜにで} \ja{はらいます。}} — \emph{kozeni de haraimasu} — \textquotedblleft{}I pay with small coins.\textquotedblright{}
\end{quote}

\begin{cousinweb}
\textbf{\ja{こぜに}} (\emph{kozeni}) is built from native \emph{ko-}, \textquotedblleft{}small\textquotedblright{}, and \emph{zeni}, \textquotedblleft{}coin\textquotedblright{}; KANJIDIC2 lists \emph{zeni} as a kun reading, but it looks like an early loan from Chinese (Middle Chinese \emph{dzjen}), so its history is not settled (JMdict, KANJIDIC2).
\end{cousinweb}

\begin{grammarlens}[title={What you've built}]
One more word for everyday life at A2.
\end{grammarlens}

\subsection*{Your turn}
\begin{itemize}
\raggedright
  \item \emph{Say it:} \emph{kozeni}
  \item \emph{Say it:} \emph{kozeni}, once more
  \item \emph{Say it:} \emph{kozeni de haraimasu} — I pay with small coins.
  \item \emph{Recall:} say the Japanese for cash, then the Japanese for banknote, then say \emph{kozeni} again
\end{itemize}

\begin{tcolorbox}[breakable,colback=teal!4,colframe=teal!35!black,title={Before you move on}]
What does \ja{こぜに} mean? (\textquotedblleft{}Small coins; small change\textquotedblright{}.) Say it once more.
\end{tcolorbox}
